% a07_k01steps_a2: part of app:a07_k01steps (input from sec/a07_k01steps.tex): cut data of per-label full classes and the use of the two regrouping lemmas.

\begin{lemma}[cut data of a per-label full class]\label{lem:a07:full}
Let $\cG$ be a per-label full class of type $T$, order $o$ and label multiset $\cL$, and group the unoriented cut terms by their \emph{cut datum} $\delta$: the components $K$ of $\Gamma\setminus S$ with their label sets $\cL_K$ and their orders, at each $K$ the numbers of $S$-ends of each kind and of original external fields of each kind, and the incidence of the $S$-objects (which components each $S$-line joins, in which component each $S$-leg sits). Then:
\begin{enumerate}[label=\textup{(\alph*)},nosep,leftmargin=*]
\item Label the $S$-objects and fix one matching of their ends. The cut terms with datum $\delta$ correspond, $s_\delta$-to-one, to the tuples $\prod_K\cK_K$, where $\cK_K$ is the per-label full class of the component (all connected graphs with label multiset $\cL_K$ and the order of $K$ whose unpaired legs carry the prescribed fields, $S$-end kernels and original fields, in all assignments, with their counterterm members, \emph{members without internal line and of order $0$ included}), and $s_\delta$ is the order of the group of relabellings of identical $S$-objects of $\delta$, the exchange of the two ends of an $S$-line with both ends in one component included.
\item At fixed labels and separable indices, $\cR^{(\delta)}=(-1)^{|S|}c_*s_\delta^{-1}\prod_KE(\cK_K;\text{fields})$.
\item Every orbit of \cref{lem:a07:orb} lies in one group; every group is closed under the exchange of the two scalar legs of every two-scalar-leg or $\phi\phi A$ subgraph, the component itself included, and under the completion of a one-vector-leg subgraph to the two-vector-leg class of its first background-Taylor coefficient. Hence each group is a union of orbits also for the moves enlarged by these two.
\item $\sum_S\cR^{(S)}=\sum_\delta\cR^{(\delta)}$.
\end{enumerate}
\end{lemma}

\begin{proof}
(a) Cutting lands in $\prod_K\cK_K$: each component is connected, has the label set $\cL_K$ and the prescribed unpaired legs, and every counterterm vertex of $\Gamma$ lies in one component with labels inside $\cL_K$. Conversely, given members $K_i\in\cK_K$, an assignment of the labelled stub fields to their unpaired legs (included in $\cK_K$) and the fixed matching, gluing the $S$-lines gives a graph with label multiset $\bigcup_K\cL_K=\cL$, the uncut type $T$ and the order $o$, connected because the incidence multigraph of $\delta$ is connected (the original $\Gamma$ was); it is therefore a member of the full class $\cG$, with $S$ among its lines and legs. A two-scalar-leg subgraph of $K_i$ avoids $S$, so by \cref{lem:a07:orb}(a) it is a two-scalar-leg subgraph of $\Gamma$, and its counterterm member corresponds to $\Gamma$'s; conversely a counterterm vertex of $\Gamma$ is a vertex of one component. Two labelled tuples give the same unlabelled cut term if and only if they differ by a relabelling of identical $S$-objects, and such a relabelling maps the product of full classes to itself (each $\cK_K$ contains all assignments of its stub fields); so the map is $s_\delta$-to-one. (b) As in \cref{lem:a07:orb}(d). (c) The moves keep $\delta$: labels by (A4), stubs by \cref{lem:a07:orb}(f), incidence because $S$-lines survive. A full class contains every pairing of its legs, and the two further moves are re-pairings inside a component with the same labels and stubs (the exchange is a bijection of Wick pairings; the first background-Taylor coefficient of the full connected generating function of a one-vector-leg type is the full two-vector-leg class, since the generating function depends on the background and on the other vector legs only through their sum). (d) The groups partition the unoriented cut terms.
\end{proof}

\begin{remark}[how the two lemmas are used]\label{rem:a07:use}
When a renormalisation bound (\cite[Prop.~1, (1.33)]{Ba3} or its cube form) is applied to a telescoped family, it is applied per orbit: by \cref{lem:a07:orb}(c), (d) each orbit is a product of complete classes of the components, each carrying at least one external field (an $S$-end or an original field), with the cubes of the original graph, so the bound applies to each component with its tree decay, and the components are joined along a spanning tree of the cut graph's component multigraph by the decay of the cut kernels (\cite[(3.2), p.~432]{Ba3}: ``the exponential factor in (1.33) is obtained from the estimates of external fields''). The gain of a cut kernel is read at the maximum over the end cubes of the orbit, which is legitimate because the kernel bounds hold at every pair of cubes at the required depth (\cref{lem:a07:K3}); where a consumer needs the gain at one member's cubes, any two cubes of the fixed vertex set $W$ are within $d(W)$, so the conversion costs $e^{2|S|\delta'_\sigma d(W)}$ with the gain taken at the reduced rate $\delta'_\sigma=\delta'/(4|S|)$, which half of the tree factor $e^{-\delta'd(W)}$ absorbs. At zero background the orbits of (P), (W), (X) suffice. At a nonzero regular background the renormalisation bound is proved for per-label full classes (the touched two-scalar-leg subgraph becomes a $\phi\phi A$ node whose partner is its mirror, a move not among (P), (W), (X)); there the regrouping is by the groups of \cref{lem:a07:full}, and the members of the component classes without internal line or of order $0$ are treated as in the proof of \cref{lem:a07:S}.
\end{remark}

